%% GENERATED by python/paper/build.py from results/ -- do not edit by hand.
%% Rebuild:  .venv\Scripts\python.exe python\paper\build.py --only UK_ESC_IncomeDistr
%% Source:   results/esc/escExperimentsUK.csv
\begin{table}[!htb]
\centering
\begin{threeparttable}
\caption{Endogenous design and the French income distribution in the UK}
\label{table:UK_ESC:incomeDistr}
\renewcommand{\arraystretch}{1.25}
\begin{tabularx}{\textwidth}{p{2.6cm}YYYYY}
\toprule
\textbf{Scenario} & \textbf{CRRA} ($\rho$) & $\bm{\theta}$ \textbf{(2020)} & \textbf{Tax rate} & \textbf{Savings rate} & \textbf{Avg. workweek} \\
\midrule 
 & 0.5 & 0.56 & 11.86\% & 19.11\% & 34.73 \\
Baseline & 1.0 & 0.56 & 11.86\% & 17.25\% & 34.75 \\
 & 2.0 & 0.56 & 11.86\% & 15.81\% & 34.74 \\[.5em]\hline\\[-.75em]
 & 0.5 & 0.56 & 11.79\% & $+0.11$ p.p. & 34.63 \\
Exogenous $\theta$ & 1.0 & 0.56 & 12.03\% & $-0.03$ p.p. & 34.58 \\
 & 2.0 & 0.56 & 12.14\% & $-0.06$ p.p. & 34.56 \\[.5em]\hline\\[-.75em]
 & 0.5 & 0.57 & 11.81\% & $+0.03$ p.p. & 34.63 \\
Endogenous $\theta$ & 1.0 & 0.55 & 12.02\% & $-0.04$ p.p. & 34.59 \\
 & 2.0 & 0.50 & 12.12\% & $-0.03$ p.p. & 34.53 \\[.5em]\hline\\[-.75em]
 & 0.5 & 1.00 & 21.29\% & $-4.45$ p.p. & 35.24 \\
France & 1.0 & 1.00 & 21.29\% & $-3.81$ p.p. & 35.24 \\
 & 2.0 & 1.00 & 21.29\% & $-3.34$ p.p. & 35.24 \\[.5em]\hline\\[-.75em]
\bottomrule
\end{tabularx}
\begin{tablenotes}[flushleft]
\footnotesize
\item[] \textit{Note:} The cost specification and the units are those of table \ref{table:US_ESC:ageing}, at the UK\textquotesingle s own cost parameter, calibrated per $\rho$ so that the UK electorate re-elects its observed design $\theta^{\ast} = 0.560$: $\lambda = 15.089, 7.264, 2.786$ at $\rho = 0.5, 1.0, 2.0$, with $\beta$ imposed from the US and $\omega$ recalibrated at each trial value. France\textquotesingle s income groups are cut at the UK\textquotesingle s own income percentiles (table \ref{table:a_US:CalibFRUK}). The France row is France's own calibrated path rather than a counterfactual on the UK model: it carries France\textquotesingle s own $\omega$ as well as its characteristics, and its workweek is a calibration target rather than a prediction.
\end{tablenotes}
\end{threeparttable}
\end{table}
